\begin{problem}[1]
  Let $J = \begin{bmatrix}
    0 & 1 \\
    -1 & 0 \\
  \end{bmatrix}$. Compute $e^{tJ}$ (where $t \in \R$).
\end{problem}

\begin{solution}
  Let $A = tJ$. The matrix exponential $e^A$ is computed using power series, defined by:
  \begin{equation}\label{eq:matrix_exponential}
    e^a = \sum_{k=0}^\infty \frac{A^k}{k!} = \frac{A^0}{0!} + \frac{A^1}{1!} + \frac{A^2}{2!} + \cdots
  \end{equation}
  Computing the powers of $A$, we have:
  \[%
    A^0 = \begin{bmatrix}
      0 & t \\
      -t & 0 \\
    \end{bmatrix}, \quad
    A^2 = \begin{bmatrix}
      -t & 0 \\
      0 & -t \\
    \end{bmatrix}, \quad
    A^3 = \begin{bmatrix}
      0 & -t \\
      t & 0 \\
    \end{bmatrix}, \quad
    A^4 = \begin{bmatrix}
      t & 0 \\
      0 & t \\
    \end{bmatrix}, \quad
    A^5 = \begin{bmatrix}
      0 & t \\
      -t & 0 \\
    \end{bmatrix} = A
  .\]%
  Notice that we repeat after $k = 4$, i.e., $A^7 = A^3$. Plugging all this into Equation~\eqref{eq:matrix_exponential}, we get:
  \begin{align*}
    e^A &= I + A + \frac{A^2}{2!} + \cdots \\
        &= I + J - \frac{t^2}{2!}I - \frac{t^3}{3!}J + \frac{t^4}{4!}I + \frac{t^5}{5!}J - \frac{t^6}{6!}I + \frac{t^7}{7!}J + \cdots \\
        &= \left(1 - \frac{t^2}{2!} + \frac{t^4}{4!} - \frac{t^6}{6!} + \cdots\right)I + \left(t - \frac{t^3}{3!} + \frac{t^5}{5!} - \frac{t^7}{7!} + \cdots\right)J
  .\end{align*}
  Notice that:
  \[%
    1 - \frac{t^2}{2!} + \frac{t^4}{4!} - \frac{t^6}{6!} + \cdots \cos(t) \aand t - \frac{t^3}{3!} + \frac{t^5}{5!} - \frac{t^7}{7!} + \cdots = \sin(t)
  .\]%
  Thus, we have:
  \[%
    e^A = \cos(t)I + \sin(t)J = \begin{bmatrix}
      \cos(t) & \sin(t) \\
      -\sin(t) & \cos(t) \\
    \end{bmatrix}
  .\qedhere\]%
\end{solution}

\begin{problem}[2]
  Recall the Bloch equation for the magnetization vector $M : \R \times \R^3 \to \R^3$ in the external magnetic field $B = B(t, x, y, z) \in \R^3$:
  \[%
    \odv{}{t} M = \gamma[M \times B] = -\frac{1}{T_2}M^{\perp} - \frac{1}{T_1}[M^\parallel - M_0]
  ,\]%
  where $T_1$, $T_2$ are the relaxation constants, $M^\perp$ is the transverse component of $M$ relative to $B$, and $M^\parallel$ is the longitudinal component of $M$ relative to $B$.

  \noindent Let $b_0 \in \R$, $\omega_0 := \gamma b_0$ and define
  \[%
    U(t) = \begin{bmatrix}
      \cos(\omega_0t) & \sin(\omega_t) & 0 \\
      -\sin(\omega_0t) & \cos(\omega_t) & 0 \\
      0 & 0 & 0 \\
    \end{bmatrix}
  .\]%
  Define $m = m(t, x, y, z)$ via
  \[%
    M(t, x, y, z) = U(t)m(t, x, y, z)
  .\]%
  Show that $m$ satisfies the equation
  \[%
    \odv{}{t} m = \gamma[m \times B_{\text{eff}}] - \frac{1}{T_2}m^\perp - \frac{1}{T_1}[m^\parallel - M_0]
  ,\]%
  where
  \[%
    B_{\text{eff}} = U^{-1}(t)B - (0, 0, b_0)
  ,\]%
  and $m^\perp$, $m^\parallel$ are defined analogously.

  \noindent \textit{Suggestion.} First set $T_1 = T_2 = \infty$ (i.e. ignore the relaxation terms) and do the calculation. Then, check the relaxation terms once you've got the hang of it.
\end{problem}

\begin{solution}
  We first consider the case without relaxation ($T_1, T_2 \to \infty$). Given $M = Um$, we differentiate with respect to $t$:
  \[%
    \odv{M}{t} = \odv{U}{t}m + U\odv{m}{t}
  .\]%
  Substituting the Bloch equation $\odv{M}/{t} = \gamma(M \times B)$, we have:
  \[%
    \gamma(Um \times B) = \odv{U}{t}m + U\odv{m}{t}
  .\]%
  To isolate $\odv{m}/{t}$, multiply by $U^{-1} = U^T$:
  \[%
    \odv{m}{t} = U^{-1}\gamma(Um \times B) - U^{-1}\odv{U}{t}m
  .\]%
  Using the property $U^{-1}(Ra \times Rb) = a \times U^{-1}b$ for rotation matrices, the first term becomes $\gamma(m \times U^{-1}B)$. Now, we compute $U^{-1}\odv{U}/{t}$:
  \[%
    \odv{U}{t} = \omega_0 \begin{bmatrix}
      -\sin(\omega_0t) & \cos(\omega_0t) & 0 \\
      -\cos(\omega_0t) & -\sin(\omega_0t) & 0 \\
      0 & 0 & 0 \\
    \end{bmatrix} \implies U^{-1}\odv{U}{t} = \begin{bmatrix}
      0 & \omega_0 & 0 \\
      -\omega_0 & 0 & 0 \\
      0 & 0 & 0 \\
    \end{bmatrix}
  .\]%
  Note that for any vector $m$, $(U^{-1}\odv{U}/{t})m = \vec{\omega} \times m$, where $\vec{\omega} = (0, 0, -\omega_0)$. Thus:
  \[%
    -U^{-1}\odv{U}{t}m = m \times (0, 0, -\omega_0) = \gamma(m \times (0, 0, -b_0))
  .\]%
  Combining these, we get:
  \[%
    \odv{m}{t} = \gamma(m \times U^{-1}B) + \gamma(m \times (0, 0, -b_0)) = \gamma [m \times (U^{-1}B - (0, 0, b_0))]
  .\]%
  This yields $B_{\text{eff}} = U^{-1}B - (0, 0, b_0)$. Finally, since $U$ is a rotation about the $z$-axis (the longitudinal axis), it preserves the decomposition of $M$ into $M^\parallel$ and $M^\perp$. Specifically, $U^{-1}M^\parallel = m^\parallel$ and $U^{-1}M^\perp = m^\perp$, ensuring the relaxation terms transform linearly into the same functional form for $m$.
\end{solution}

\begin{problem}[3]
  Show that if $A$ is skew-symmetric (i.e. $A^* = -A$), then $e^A$ is unitary (i.e. $(e^A)^* (e^A) = I$).
\end{problem}

\begin{solution}
  The conjugate transpose of a matrix exponential satisfies $(e^A)^* = e^{A^*}$. Substituting the skew-symmetry condition $A^* = -A$, we obtain:
  \[%
    (e^A)^* = e^{-A}
  .\]%
  To check if $e^A$ is unitary, we compute the product:
  \[%
    (e^A)^* (e^A) = e^{-A} e^A
  .\]%
  Since $-A$ and $A$ commute (specifically, $-A A = A (-A) = -A^2$), we can use the property of matrix exponentials $e^X e^Y = e^{X+Y}$ to write:
  \[%
    e^{-A} e^A = e^{-A+A} = e^0 = I
  .\]%
  Thus, $e^A$ is unitary.
\end{solution}

\begin{problem}[4]
  Identify $\R^2$ with $\C$ via $(z_1, z_2) \leftrightarrow z_1 + iz_2$. Let $z = z_1 + iz_2$. Show that the linear transformation $T_z : \C \to \C$ given by $T_z(w) = zw$ is represented (as a linear transformation from $\R^2$ to $\R^2$) by the matrix
  \[%
    \begin{bmatrix}
      z_1 & -z_2 \\
      z_2 & z_1 \\
    \end{bmatrix}
  .\]%
\end{problem}

\begin{solution}
  To find the matrix representation of a linear transformation $T$ from $\R^2$ to $\R^2$, we examine the images of the standard basis vectors $\e_1 = (1, 0)$ and $\e_2 = (0, 1)$. In the complex identification, $\e_1$ corresponds to $w = 1 + 0i = 1$ and $\e_2$ corresponds to $w = 0 + 1i = ij$. Now, we apply the transformation $T_z(w)$. For $w = 1$:
  \[%
    T_z(1) = z \cdot 1 = z_1 + iz_2
  .\]%
  In $\R^2$, this corresponds to the vector $\begin{bmatrix} z_1 \\ z_2 \end{bmatrix}$. This forms the first column of our matrix. For $w = 2$:
  \[%
    T_z(i) = z \cdot i = (z_1 + iz_2)i = z_1 i + i^2 z_2 = -z_2 + iz_1 
  .\]%
  In $\R^2$, this corresponds to the vector $\begin{bmatrix} -z_2 \\ z_1 \end{bmatrix}$. This forms the second column of our matrix. Combining these columns, we obtain the matrix representation:
  \[%
    [T_z] = \begin{bmatrix}
      z_1 & -z_2 \\
      z_2 & z_1 \\
    \end{bmatrix}
  ,\]%
  which matches the required form.
\end{solution}

\begin{problem}[5]
  Suppose $B \in \R^3$ and $M = M(t) \in \R^3$ satisfies
  \[%
    \odv{}{t} M = M \times B
  .\]%
  Show that $|M(t)|^2$ is constant.
\end{problem}

\begin{solution}
  Recall that the squared magnitude of a vector can be expressed using the dot product: $|M(t)|^2 = M \cdot M$. To show this is constant, we differentiate with respect to $t$:
  \[%
    \odv{}{t} |M(t)|^2 = \odv{}{t} (M \cdot M)
  .\]%
  Using the product rule for dot products, we have:
  \[%
    \odv{}{t} (M \cdot M) = \odv{M}{t} \cdot M + M \cdot \odv{M}{t} = 2 M \cdot \odv{M}{t}
  .\]%
  Now, substitute the given differential equation $\odv{M}/{t} = M \times B$:
  \[%
    \odv{}{t} |M(t)|^2 = 2 M \cdot (M \times B)
  .\]%
  By the properties of the scalar triple product, the vector $M \times B$ is orthogonal to both $M$ and $B$. Specifically, $M \cdot (M \times B) = 0$ for any vector $M$. Therefore:
  \[%
    \odv{}{t} |M(t)|^2 = 2(0) = 0
  .\]%
  Since the derivative is zero for all $t$, $|M(t)|^2$ must be constant.
\end{solution}
